% OpenStax Calculus Volume 2 -- Section 4.4 Exercises: The Logistic Equation
% Exercises reproduced from OpenStax, licensed under CC BY 4.0.
% Source: https://openstax.org/books/calculus-volume-2/pages/4-4-the-logistic-equation
\documentclass{exercises}
\title{OpenStax Calculus Volume 2}
\subtitle{Section 4.4 Exercises: The Logistic Equation}
\sourceurl{https://openstax.org/books/calculus-volume-2/pages/4-4-the-logistic-equation}
\author{}
\date{}

\begin{document}
\maketitle


For the following problems, consider the logistic equation in the form $P'=CP-P^{2}$. Draw the directional field and find the stability of the equilibria.

\begin{enumerate}[start=1]
\item $C=3$
\item $C=0$
\item $C=-3$
\item Solve the logistic equation for $C=10$ and an initial condition of $P(0)=2$.
\item Solve the logistic equation for $C=-10$ and an initial condition of $P(0)=2$.
\item A population of deer inside a park has a carrying capacity of $200$ and a growth rate of $2\%$. If the initial population is $50$ deer, what is the population of deer at any given time?
\item A population of frogs in a pond has a growth rate of $5\%$. If the initial population is $1000$ frogs and the carrying capacity is $6000$, what is the population of frogs at any given time?
\item \techrequired Bacteria grow at a rate of $20\%$ per hour in a petri dish. If there is initially one bacterium and a carrying capacity of $1$ million cells, how long does it take to reach $500{,}000$ cells?
\item \techrequired Rabbits in a park have an initial population of $10$ and grow at a rate of $4\%$ per year. If the carrying capacity is $500$, at what time does the population reach $100$ rabbits?
\item \techrequired Two monkeys are placed on an island. After $5$ years, there are $8$ monkeys, and the estimated carrying capacity is $25$ monkeys. When does the population of monkeys reach $16$ monkeys?
\item \techrequired A butterfly sanctuary is built that can hold $2000$ butterflies, and $400$ butterflies are initially moved in. If after $2$ months there are now $800$ butterflies, when does the population get to $1500$ butterflies?
\end{enumerate}

The following problems consider the logistic equation with an added term $f$ for depletion, either through death or emigration:

\[
P'=0.4P\left(1-\frac{P}{10000}\right)-f.
\]

\begin{enumerate}[resume]
\item \techrequired The population of trout in a pond is given by $P'=0.4P(1-\frac{P}{10000})-400$, where $400$ trout are caught per year. Use your calculator or computer software to draw a directional field and draw a few sample solutions. What do you expect for the behavior?
\item In the preceding problem, what are the stabilities of the equilibria $0<P_{1}<P_{2}$?
\item \techrequired For the preceding problem, use software to generate a directional field for the value $f=400$. What are the stabilities of the equilibria?
\item \techrequired For the preceding problems, use software to generate a directional field for the value $f=600$. What are the stabilities of the equilibria?
\item \techrequired For the preceding problems, consider the case where a certain number of fish are added to the pond, or $f=-200$. What are the nonnegative equilibria and their stabilities?
\end{enumerate}

It is more likely that the amount of fishing is governed by the current number of fish present, so instead of a constant number of fish being caught, the rate is proportional to the current number of fish present, with proportionality constant $k$, as


$P'=0.4P(1-\frac{P}{10000})-kP$.

\begin{enumerate}[resume]
\item \techrequired For the previous fishing problem, draw a directional field assuming $k=0.1$. Draw some solutions that exhibit this behavior. What are the equilibria and what are their stabilities?
\item \techrequired Use software or a calculator to draw directional fields for $k=0.4$. What are the nonnegative equilibria and their stabilities?
\item \techrequired Use software or a calculator to draw directional fields for $k=0.6$. What are the equilibria and their stabilities?
\item Solve this equation, assuming a value of $k=0.05$ and an initial condition of $2000$ fish.
\item Solve this equation, assuming a value of $k=0.05$ and an initial condition of $5000$ fish.
\end{enumerate}

The following problems add in a minimal threshold value for the species to survive, $T$, which changes the differential equation to $P'(t)=rP(1-\frac{P}{K})(1-\frac{T}{P})$.

\begin{enumerate}[resume]
\item Draw the directional field of the threshold logistic equation, assuming $K=10$, $r=0.1$, and $T=2$. When does the population survive? When does it go extinct?
\item For the preceding problem, solve the logistic threshold equation, assuming the initial condition $P(0)=P_{0}$.
\item Bengal tigers in a conservation park have a carrying capacity of $100$ and need a minimum of $10$ to survive. If they grow in population at a rate of $1\%$ per year, with an initial population of $15$ tigers, solve for the number of tigers present.
\item A forest containing ring-tailed lemurs in Madagascar has the potential to support $5000$ individuals, and the lemur population grows at a rate of $5\%$ per year. A minimum of $500$ individuals is needed for the lemurs to survive. Given an initial population of $600$ lemurs, solve for the population of lemurs.
\item The population of Florida panthers in the southwestern tip of Florida has an estimated carrying capacity of $250$ and grows at a rate of $0.25\%$ per year, and there must be $25$ for the population to survive. With an initial population of $30$ Florida panthers, how many years will it take to get the Florida panthers off the endangered species list (at least $100$)?
\end{enumerate}

The following questions consider the Gompertz equation, a modification for logistic growth, which is often used for modeling cancer growth, specifically the number of tumor cells.

\begin{enumerate}[resume]
\item The Gompertz equation is given by $P'(t)=\alpha \ln(\frac{K}{P(t)})P(t)$. Draw the directional fields for this equation assuming all parameters are positive, and given that $K=1$.
\item Assume that for a population, $K=1000$ and $\alpha=0.05$. Draw the directional field associated with this differential equation and draw a few solutions. What is the behavior of the population?
\item Solve the Gompertz equation for generic $\alpha$ and $K$ and $P(0)=P_{0}$.
\item \techrequired The Gompertz equation has been used to model tumor growth in the human body. Starting from one tumor cell on day $1$ and assuming $\alpha=0.1$ and a carrying capacity of $10$ million cells, how long does it take to reach ``detection'' stage at $5$ million cells?
\item \techrequired It is estimated that the world human population reached $3$ billion people in $1959$ and $6$ billion in $1999$. Assuming a carrying capacity of $16$ billion humans, write and solve the differential equation for logistic growth, and determine what year the population reached $7$ billion.
\item \techrequired It is estimated that the world human population reached $3$ billion people in $1959$ and $6$ billion in $1999$. Assuming a carrying capacity of $16$ billion humans, write and solve the differential equation for Gompertz growth, and determine what year the population reached $7$ billion. Was logistic growth or Gompertz growth more accurate, considering world population reached $7$ billion on October 31, 2011?
\item Show that the population grows fastest when it reaches half the carrying capacity for the logistic equation $P'=rP(1-\frac{P}{K})$.
\item When does population increase the fastest in the threshold logistic equation $P'(t)=rP(1-\frac{P}{K})(1-\frac{T}{P})$?
\item When does population increase the fastest for the Gompertz equation $P'(t)=\alpha \ln(\frac{K}{P(t)})P(t)$?
\end{enumerate}

Below is a table of the populations of whooping cranes in the wild from $1940$ to $2000$. The population rebounded from near extinction after conservation efforts began. The following problems consider applying population models to fit the data. Assume a carrying capacity of $10{,}000$ cranes. Fit the data assuming years since $1940$ (so your initial population at time $0$ would be $22$ cranes).

\begin{enumerate}[resume]
\item Find the equation and parameter $r$ that best fit the data for the logistic equation.
\item Find the equation and parameters $r$ and $T$ that best fit the data for the threshold logistic equation.
\item Find the equation and parameter $\alpha$ that best fit the data for the Gompertz equation.
\item Graph all three solutions and the data on the same graph. Which model appears to be most accurate?
\item Using the three equations found in the previous problems, estimate the population in $2010$ (year $70$ after conservation). The real population measured at that time was $437$. Which model is most accurate?
\end{enumerate}

\end{document}
